\section{Affective Precision Update and Policy Selection}
\label{app:precision_update}

This appendix gives the $\beta$-update rule, policy-precision mapping, and partner-choice policy-selection procedure.

\subsection{Categorical beta support}

Each partner maintains a categorical posterior $q(\beta_k)$ over discrete levels $\{0.5, 0.67, 1.0, 1.5, 2.0\}$, updated each round from affective charge $\phi_k$. The posterior is categorical rather than continuous; this keeps the inverse-$\beta$ convention identifiable while preserving a compact implementation of the precision dynamics.

\subsection{Persistence prior}

The update proceeds in three steps. First, a persistence prior smooths the previous posterior via a tridiagonal transition matrix:
\begin{equation}
p(\beta_k) \leftarrow T \cdot q_{\text{prev}}(\beta_k),
\label{eq:persistence}
\end{equation}
where $T$ is tridiagonal with persistence parameter 0.8 and reflecting boundaries, encoding the prior expectation that precision evolves slowly.

\subsection{Charge-based pseudo-likelihood}

Second, a pseudo-likelihood is constructed directly from affective charge, mapping positive charge to high precision (low $\beta$) and negative charge to low precision (high $\beta$):
\begin{equation}
\log \ell(\beta_\ell) = \phi_k \cdot \frac{1}{\beta_\ell},
\label{eq:pseudo_lik}
\end{equation}
where $\beta_\ell$ ranges over the five support points. This is a hand-designed update rule (not derived from the POMDP generative model) chosen to satisfy monotonicity: strong positive charge reinforces low $\beta$, strong negative charge reinforces high $\beta$, and the effect scales with effective precision at each level.

\subsection{Posterior update and policy precision}

Third, the posterior is normalized:
\begin{equation}
q(\beta_k) \propto \exp(\log \ell(\beta_\ell)) \cdot p(\beta_k).
\label{eq:posterior_update}
\end{equation}
Policy precision is then set from the posterior mean as given in Eq.~3 of the main text. The $\beta_k$ posterior is maintained as an external auxiliary tracker rather than a proper hidden state in the generative model, which simplifies implementation at the cost of losing the agent's ability to form prior expectations over its own future affective states (discussed in Section~4).

\subsection{Centered cross-partner policy logits}

In partner-choice regimes, precision must sharpen policies within a partner branch without creating artifacts when branches are compared across partners. The centered logit procedure is given in Eq.~4 of the main text. The agent samples or selects policies from the combined set of centered logits across all partner branches. Centering preserves the mean score of each partner branch while allowing $\gamma_k$ to sharpen or flatten the within-partner policy distribution, preventing low-precision scaling from making uniformly poor partner branches spuriously attractive in cross-partner comparison.

\subsection{Ablation variants}

The four affect-runtime variants are defined in Section~2.5. For completeness: the no-affect control disables the $\beta$ tracker entirely and fixes $\gamma_k = \gamma_{\text{base}}$; the tracked-only ablation updates $q(\beta_k)$ normally but fixes $\gamma_k = \gamma_{\text{base}}$, allowing the precision estimate to be analyzed without influencing action selection; the default partner-local variant applies the full update pipeline described in B.1--B.4 and sets $\gamma_k$ from each partner's posterior mean; and the shared-$\beta$ ablation maintains a single pooled $\beta$ state updated from all partner evidence while keeping per-partner POMDP beliefs intact, with all partner branches inheriting precision from the shared expected $\beta$.
